
\section{HCP} \label{s:hcp}
Hitachi Content Platform (HCP) is a multipurpose distributed \os{} designed to support large-scale repositories of unstructured data. Physically, HCP is a collection of storage servers (nodes), referred to as a \emph{cluster}, that use a private back-end network for inter-node communications and a public front-end network for client communication. A gossip protocol is used determine active cluster membership. In the event of temporal failure of any node or service requests are automatically vectored to active nodes for fulfillment.

The system is logically divided into a collection of \emph{tenants} and \emph{namespaces}; tenants are the administrative unit, namespaces are the storage unit. Tenants are critical for cloud / service provider deployments where it's important not only to logically separate client data, but also to divide the administration of each virtual instance of HCP. The system is designed so that no one administrator has dominion over the system as a whole, but only over their assigned tenant. A tenant administrator can create a collection of namespaces that will hold user data.

The value proposition of a multitenant system is largely the same as any shared pool of physical resources such as SANs: Sharing reduces cost by pooling physical resources. In the process they also expose data on shared storage by unauthorized users and overwrites by multiple clients \cite{yoshida1999lun}. To counter this HCP employs a complex of security measures. These include strictly dividing administrative functions, defined user access rights and restrictions, access controls granular to individual objects, and all data is optionally encrypted, both in-flight and on disk.

Client access to the system is via a number of software gateways, distinguished by protocol. Presently the following protocols are supported: HTTP (REST), S3, CIFS, NFS, SMTP and WebDAV. Objects that were ingested using any protocol are immediately accessible through any other supported protocol. These protocols can be used to access the data with a web browser, HCP client tools, third-party applications, Windows Explorer, or native Windows or Unix tools \cite{MR:2013:HCP-ARCH}.

HCP provides high availability with strong consistency guarantees. Replication comes in two forms: intra- and inter-cluster. Intra-cluster replication is synchronous to the write path to ensure that all data ingested is successfully persisted to disk before returning acknowledgement to the client. Consistent hashing is used to distribute data among the nodes. Within the strictures of assuring replicas are assigned separate fault domains, writes bias toward lightly loaded nodes. Inter-cluster replication is asynchronous to the write path to ensure low write latency while providing geographic dispersion of objects. HCP uses XML\footnote{{\href{http://www.w3.org/XML}{http://www.w3.org/XML}}} as the serialization format for persisting \cmd{} and XPATH\footnote{\href{http://www.w3.org/TR/xpath}{http://www.w3.org/TR/xpath}} structured queries against that \md{}. This allows the system to return \emph{specific} answers to user queries, rather than a collection of \emph{potential} matches.

HCP employs redundancy at multiple levels: the object reference database, user data (data objects), and system and custom (user) \md{}. Likewise, all hardware components are redundant so there is no single point of failure. Commonly the focus for data protection centers on protecting user data. However, for database systems it's actually the protection of the index that's most important, as a loss of the pointer to the data is tantamount to the loss of the data itself. HCP shards the object reference database amongst distinct fault domains, both for performance and protection. Further, in the event of catastrophic failure, a \emph{scavenging service} is employed to reconstruct the \db{} from the constituent \md{} persisted to disk as part of the write path.

The data model is one of immutable data objects with mutable \md{}. To simulate in-place updates of data objects HCP supports object versioning, i.e. the capability of a namespace to create, store, and manage multiple versions of objects within the repository. Capacity efficiency is achieved through a combination of compression and duplicate elimination.

In addition to the actual data persistence component, the system needs to have scalable and robust solutions for load balancing, membership and failure detection, content verification, version control, retention policies, disposition of objects no longer under retention, compression, encryption, failure recovery, replica synchronization, overload handling, state transfer, garbage collection, concurrency and job scheduling, request marshaling, request routing, system monitoring and alarming, and configuration management.

The result is a system that provides data management as a service, referred to as \emph{Storage as a Service}, in a manner that makes both client development and management of the system simple. The design goal is that neither clients nor system administrators should be aware of the detail of the multiple software services needed to keep a multi-petabyte repository housing billions of objects coherent and responsive to client requests.
